\section{Fixed downstream substrate, target class, and exact conditional}
\label{sec:targetclass}

The present manuscript is not foundational in the sense of re-deriving admissibility, standing,
reference, or irreversibility from below. It is downstream. The first task is therefore to state
explicitly the fixed-domain structural package on which the measurement argument runs.

\subsection{Fixed downstream substrate}

\begin{definition}[Fixed downstream substrate]
\label{def:fixed-downstream-substrate}
A \emph{fixed downstream substrate} for the present manuscript is a same-scope, admissibility-governed
regime in which the following are already in force.
\begin{enumerate}
    \item Admissibility is a forced governance boundary for the target class of non-degenerate,
    proposition-bearing reasoning.
    \item Standing is bivalent on the evaluated fragment and fail-closed when evaluation is not
    instantiated.
    \item Same-scope admissibility classification is unique: there is no second admissibility
    predicate, no admissibility-relevant boundary selector, and no admissibility-relevant
    parameterization on the fixed scope.
    \item The admissible interior is unique up to lawful equivalence.
    \item Standing is conserved: there is no same-scope repair, no generator, no primitive carrier,
    and no same-scope transfer without fresh target gating.
    \item Any second same-scope invariant either collapses to bookkeeping or constitutes a scope
    change.
\end{enumerate}
\end{definition}

\begin{remark}[Provenance of the substrate]
\label{rem:substrate-provenance}
Definition~\ref{def:fixed-downstream-substrate} is not a hidden premise. It is the explicit fixed-domain
package proved in the companion closure papers on kernel necessity, interface forcing, trajectory
closure, uniqueness, and structural singleton results; see
\citep{maley_foundational2026,maley_necessity2026,maley_trajectory2026,maley_unus2026,
maley_clausura2026,maley_ats2026,maley_ametric2026,maley_structura2026}. The present paper uses
that package as background and does not attempt to re-prove it in full. What it does add locally is a
compressed bridge theorem showing why, once the target explanatory burden is fixed, this substrate is
not optional for the measurement argument.
\end{remark}

\subsection{Compressed substrate-necessity bridge}

\begin{definition}[Minimal coherence triggers]
\label{def:minimal-coherence-triggers}
A same-scope discourse-practice satisfies the \emph{minimal coherence triggers} relevant to the
present paper iff it preserves reference under lawful redescription, treats admissible continuation as
compositionally available candidate material, carries non-trivial inferential content, and executes at
least one irreversible commitment on the fixed scope.
\end{definition}

\begin{lemma}[Stable reference is necessary for the target explanatory practice]
\label{lem:stable-reference-necessary}
If measurement records are used as theory-level evidential anchors under lawful same-scope
redescription, then lawful redescription must preserve what those records are about. Otherwise the
same evidential claim changes target under admissible use and no standing-bearing evidential role is
fixed.
\end{lemma}

\begin{proof}
A theory-level evidential claim is supposed to remain the same claim across lawful same-scope
redescription. If lawful redescription changes what the claim is about, then the evidential role of the
record is presentation-dependent rather than theory-relative. On the fail-closed substrate recorded in
Definition~\ref{def:fixed-downstream-substrate}, that is not a weaker kind of confirmation; it is loss
of determinate standing-bearing use. So stable reference is necessary.
\end{proof}

\begin{lemma}[Identity quotient necessity]
\label{lem:identity-quotient-necessary}
If same-scope measurement-record discourse is stable under lawful redescription, then admissible
reference assignments admit a standing-equivalence relation whose classes exhaust same-scope
standing-bearing content for the target discourse. Any same-scope standing-bearing presentation of
that content therefore factors through a quotient of admissible assignments.
\end{lemma}

\begin{proof}
Once lawful redescription preserves standing-bearing use, assignments that agree on every
standing-bearing value in the target discourse cannot be treated as distinct same-scope contents unless
one introduces an additional same-scope invariant. But the fixed downstream substrate excludes any
second same-scope invariant except as bookkeeping or scope change. Therefore standing-bearing
content must collapse to equivalence classes of assignments, and any same-scope standing-bearing
presentation must factor through the corresponding quotient.
\end{proof}

\begin{lemma}[Irreversibility forces bivalent gating]
\label{lem:irrev-forces-bivalent-gating}
If same-scope illicit evaluated acts cannot be repaired as the same acts, then evaluated acts must fall
under a fail-closed eligible/ineligible gate on the evaluated fragment. No third standing-relevant
status remains.
\end{lemma}

\begin{proof}
Without a gate there is no stable distinction between eligible and ineligible evaluated acts, so
irreversibility has nothing determinate to attach to. If a third standing-relevant status were admitted,
then some acts would be neither fully eligible nor fully ineligible and illegitimacy could be softened or
reclassified without a fresh gate. That is exactly the kind of same-act repair excluded by the
irreversibility condition. Hence irreversibility forces bivalent fail-closed gating.
\end{proof}

\begin{lemma}[Bivalence plus the AMetric boundary excludes parameterization]
\label{lem:bivalence-boundary-excludes-parameterization}
If the admissibility interface is bivalent on the evaluated fragment and the pre-admissible side is
AMetric, then no same-scope admissibility-relevant parameter, ranking, or selector can refine the
boundary.
\end{lemma}

\begin{proof}
A same-scope admissibility-relevant parameter would introduce more than the two bivalent statuses of
eligible and ineligible, or else it would provide a selector among allegedly admissible alternatives.
Either way the boundary would cease to be AMetric and non-parameterized. So bivalence together
with the AMetric boundary excludes same-scope admissibility-relevant parameterization.
\end{proof}

\begin{theorem}[Compressed substrate-necessity bridge]
\label{thm:compressed-substrate-necessity}
Let a same-scope discourse-practice satisfy the minimal coherence triggers of
Definition~\ref{def:minimal-coherence-triggers} and let it be used to underwrite the target
explanatory practice of Definition~\ref{def:target-practice}. Then the fixed downstream substrate of
Definition~\ref{def:fixed-downstream-substrate} is not optional: admissibility is forced, standing is
bivalent on the evaluated fragment, same-scope admissibility classification is unique, the admissible
interior is unique up to lawful equivalence, standing is conserved, and any second same-scope
invariant collapses to bookkeeping or scope change.
\end{theorem}

\begin{proof}
Lemma~\ref{lem:stable-reference-necessary} forces stable reference for the target explanatory
practice. Lemma~\ref{lem:identity-quotient-necessary} then forces quotient-level standing-bearing
content rather than representative-dependent content. Lemma~\ref{lem:irrev-forces-bivalent-gating}
forces fail-closed bivalent gating on the evaluated fragment, and
Lemma~\ref{lem:bivalence-boundary-excludes-parameterization} excludes same-scope admissibility-
relevant boundary parameterization or selection. The remaining fixed-domain consequences --- unique
same-scope admissibility classification, unique admissible interior, standing conservation, and the
collapse of any second same-scope invariant --- are exactly the downstream closure results proved in
\citep{maley_foundational2026,maley_necessity2026,maley_trajectory2026,maley_unus2026,
maley_clausura2026,maley_ats2026,maley_ametric2026,maley_structura2026}. Therefore, once the
target explanatory practice is fixed on a same scope, the substrate recorded in
Definition~\ref{def:fixed-downstream-substrate} is forced rather than optional.
\end{proof}

\begin{proposition}[Immediate consequences of the substrate]
\label{prop:substrate-consequences}
Within the fixed downstream substrate, any same-scope admissibility-bearing distinction about
measurement-record discourse must appear in one of only two places:
\begin{enumerate}
    \item in the standing classification of evaluated acts, or
    \item in the licensed continuation structure of the discourse built on those acts.
\end{enumerate}
No third same-scope locus of admissibility-bearing work remains.
\end{proposition}

\begin{proof}
By Definition~\ref{def:fixed-downstream-substrate}, same-scope admissibility classification is unique,
standing is conserved, and no second same-scope invariant survives except as bookkeeping or scope
change. Therefore any admissibility-bearing difference must either change which acts stand, or change
what same-scope continuation is licensed for standing-bearing acts. If it changes neither, then by the
same fixed-domain package it is semantic or representational surplus rather than admissibility-bearing
content. So no third locus remains.
\end{proof}

\subsection{Physically non-selective theories and the referential selection problem}

\begin{definition}[Physical non-selection]
\label{def:physical-nonselection}
A theory is \emph{physically non-selective} with respect to a measurement process if, after the
measurement interaction, the theory's physical description realizes multiple successor histories
containing stable measurement records and physically continuous observer-successors, while
supplying no physical law, state variable, or dynamically privileged trajectory that singles out one
such successor history as physically distinguished from the others.
\end{definition}

\begin{definition}[Diachronic referential selection problem]
\label{def:diachronic-selection-problem}
Given a physically non-selective theory, the \emph{diachronic referential selection problem} is the
problem whether the theory, together with its admissible interpretive use on the fixed downstream
substrate, determines which successor history is the referent of pre-measurement diachronic discourse.
\end{definition}

Definitions~\ref{def:physical-nonselection} and \ref{def:diachronic-selection-problem} separate two
questions that are often run together. Physical non-selection is a claim about what the physical theory
itself delivers: multiplicity without a physically distinguishing selector. The diachronic referential
selection problem is the further question whether that physical package, together with whatever
admissible same-scope use it supports, determines which successor history is the referent of discourse
about what the pre-measurement subject will observe, remember, or use as evidence.

Everettian quantum mechanics is the paradigm case. Unitary dynamics and decoherence can stabilize
multiple post-measurement branch-histories, each carrying a determinate branch-relative observer and
a determinate branch-relative record, while leaving the physical theory itself non-selective among
them. The issue is not whether the theory assigns numerical weights to those histories. The issue is
whether the theory itself fixes which physically realized history is the referent of the measurement
record discourse of an internal subject.

\subsection{Exact conditional and target discourse}

The manuscript runs under one standing conditional.

\begin{claimbox}
\textbf{Standing conditional.} If a quantum theory is used to underwrite truth-apt diachronic discourse
about measurement, memory, anticipation, deliberation, and evidential confirmation, then
admissible reference assignments for measurement records must either
\begin{enumerate}
    \item collapse into a single standing-bearing class on that discourse, or
    \item leave confirmation and diachronic practice assignment-relative rather than theory-relative.
\end{enumerate}
\end{claimbox}

This conditional is not a disguised collapse postulate. It says only that if a theory is assigned
explanatory work involving records as evidence, then the referents of those records cannot vary in
same-scope standing-bearing ways while the theory still claims a single theory-level evidential output.

\begin{definition}[Target explanatory practice]
\label{def:target-practice}
The \emph{target explanatory practice} of this paper is the practice in which measurement records
function as evidential anchors for theory-level claims about what was observed, what follows for
memory and anticipation, and how those observations bear on rival hypotheses.
\end{definition}

\begin{definition}[Fixed scope of the paper]
\label{def:fixed-scope-paper}
The fixed scope of the paper is the same-scope, discourse-committing use of a physically non-selective
measurement theory within the target explanatory practice of Definition~\ref{def:target-practice}.
Moves that import new ontology, new dynamics, new admissibility-relevant structure, new boundary
selectors, or new target gating count as scope changes rather than same-scope repairs.
\end{definition}

\subsection{Public vocabulary}

The main body stays close to standard foundations vocabulary while keeping the fixed-domain
background explicit.

\begin{itemize}
    \item \textbf{Measurement record}: a physically realized outcome-bearing structure used as an
    evidential anchor.
    \item \textbf{Diachronic measurement discourse}: discourse linking a pre-measurement subject to
    post-measurement records by way of memory, anticipation, deliberation, or confirmation.
    \item \textbf{Reference assignment}: a map from diachronic measurement claims to candidate
    branch-record referents.
    \item \textbf{Admissible reference assignment}: a same-scope reference assignment eligible under the
    gate-governed fixed-domain constraints stated in Section~\ref{sec:semanticcore}.
    \item \textbf{Selector}: any principle that fixes singular reference by privileging one continuation,
    history, or record as the referent relevant to the target discourse.
    \item \textbf{Standing-bearing content}: content that survives into standing classification or licensed
    continuation on the fixed downstream substrate.
    \item \textbf{Skin}: representational difference that does not survive into standing-bearing content.
\end{itemize}

\subsection{Scope of the present manuscript}

Theorems in the present manuscript do not settle every metaphysical question in quantum
foundations. Their function is narrower. On the fixed downstream substrate,
\begin{enumerate}
    \item the paper identifies the measurement problem as a fixed-domain problem of reference to
    measurement records,
    \item proves a quotient uniqueness theorem for admissible measurement-reference content,
    \item shows that the bare Everettian package is physically non-selective with respect to that
    target discourse,
    \item proves that confirmational divergence among admissible assignments is quotient-relevant and
    therefore not mere semantic skin,
    \item proves that theory-level confirmation requires fixation commutation under admissible
    redescription,
    \item proves that memory, anticipation, and agency inherit the same lineage dependence, and
    \item classifies interpretation packages by the structural locus at which singular reference is fixed,
    if it is fixed at all.
\end{enumerate}
The manuscript is therefore neither a general semantics of quantum theory nor a re-derivation of the
companion closure corpus. It is a fixed-domain measurement-reference paper built on that corpus.
